% Compile with: latexmk -xelatex report.tex
\documentclass[12pt, a4paper, oneside]{report}

\usepackage[top=2.6cm, bottom=2.6cm, left=3cm, right=2.4cm]{geometry}
% Shared preamble for report.tex and slides.tex.
% Compile both with XeLaTeX: latexmk -xelatex <file>.tex
%
% Visual identity, chosen to be distinct from the Discrete Mathematics template:
%   text     TeX Gyre Pagella (a Palatino), wider and rounder than Times
%   headings TeX Gyre Heros (a Helvetica), sans, to separate structure from prose
%   maths    TeX Gyre Pagella Math through unicode-math, so the digits inside
%            formulas match the digits in the surrounding sentence
% All three ship with TeX Live, so nothing extra has to be installed.

\usepackage{fontspec}
\usepackage{polyglossia}
\setmainlanguage{vietnamese}

\setmainfont{texgyrepagella}[
  Extension      = .otf,
  UprightFont    = *-regular,
  BoldFont       = *-bold,
  ItalicFont     = *-italic,
  BoldItalicFont = *-bolditalic,
]
\setsansfont{texgyreheros}[
  Extension      = .otf,
  UprightFont    = *-regular,
  BoldFont       = *-bold,
  ItalicFont     = *-italic,
  BoldItalicFont = *-bolditalic,
  Scale          = 0.92,
]
\setmonofont{texgyrecursor}[
  Extension      = .otf,
  UprightFont    = *-regular,
  BoldFont       = *-bold,
  Scale          = 0.88,
]

\usepackage{amsmath}
\usepackage{amssymb}
\usepackage{amsthm}
\usepackage{mathtools}
\usepackage{unicode-math}
\setmathfont{texgyrepagella-math.otf}

\usepackage{graphicx}
\usepackage{booktabs}
\usepackage{siunitx}
\usepackage{xcolor}

% Vietnamese decimal marker, so \num agrees with the decimals written by hand in
% the prose (1{,}39 rather than 1.39).
\sisetup{output-decimal-marker={,}}

% Slide copies come first on purpose: \resultgraphic must prefer them.
\graphicspath{{figures/}{../results/figures/slides/}{../results/figures/}}

% ---------------------------------------------------------------------------
% Palette. The three accents are the figure colours from src/figures.py, so a
% method named in the text can be tinted to match its curve.
% ---------------------------------------------------------------------------
\definecolor{ink}{HTML}{1A1A1A}
\definecolor{accent}{HTML}{0B5563}     % deep teal, the structural colour
\definecolor{rule}{HTML}{9AA5A8}
\definecolor{cGD}{HTML}{1F77B4}
\definecolor{cSGD}{HTML}{FF7F0E}
\definecolor{cAGD}{HTML}{2CA02C}
\definecolor{cNewton}{HTML}{D62728}

% ---------------------------------------------------------------------------
% Notation, kept identical to KE_HOACH_TRIEN_KHAI.md
% ---------------------------------------------------------------------------
\newcommand{\R}{\mathbb{R}}
\newcommand{\E}{\mathbb{E}}
\newcommand{\wstar}{w^{*}}
\newcommand{\fstar}{f^{*}}
\newcommand{\bigO}{\mathcal{O}}
\newcommand{\Batch}{\mathcal{B}}
\newcommand{\trans}{^{\top}}
\DeclareMathOperator*{\argmin}{arg\,min}
\DeclareMathOperator{\sign}{sign}
\DeclarePairedDelimiter{\norm}{\lVert}{\rVert}
\DeclarePairedDelimiter{\abs}{\lvert}{\rvert}

% ---------------------------------------------------------------------------
% Figures. A stem that has not been generated yet leaves a visible placeholder
% rather than a compilation error, so the report can be drafted while the
% experiments are still running.
% ---------------------------------------------------------------------------
\newcommand{\missingfigbox}[1]{%
  \fcolorbox{rule}{white}{\parbox[c][3.2cm][c]{0.86\linewidth}{\centering
    \ttfamily\small\detokenize{#1}.pdf\par\vspace{4pt}
    \rmfamily\small chưa được sinh ra từ notebook\par}}%
}

% The path is spelled out rather than left to \graphicspath, which lists the
% slide directory first and would otherwise hand the report the wide, reduced
% copies meant for Beamer.
%
% Placement is [tbp] with no 'h': a figure declared between two paragraphs of one
% argument would otherwise be typeset right there and cut the argument in half.
\newcommand{\resultfig}[3]{%
  \begin{figure}[tbp]
    \centering
    \IfFileExists{../results/figures/#1.pdf}
      {\includegraphics[width=0.82\linewidth]{../results/figures/#1.pdf}}
      {\missingfigbox{#1}}
    \caption{#2}
    \label{#3}
  \end{figure}%
}

% The mandatory pair, stacked rather than side by side. The source figures are
% 6 by 4 inches with 10 pt labels; at half the text width that becomes about 5 pt
% and the legends stop being readable, which defeats the purpose of printing two
% panels. Stacking costs most of a page per comparison and keeps both legible.
\newcommand{\panellabel}[1]{{\sffamily\small\color{accent}#1}}
\newcommand{\resultpair}[3]{%
  \begin{figure}[tbp]
    \centering
    \panellabel{(a) theo số vòng lặp}\par\vspace{2pt}
    \IfFileExists{../results/figures/#1_iters.pdf}
      {\includegraphics[width=0.8\linewidth]{../results/figures/#1_iters.pdf}}
      {\missingfigbox{#1_iters}}
    \par\vspace{10pt}
    \panellabel{(b) theo thời gian chạy}\par\vspace{2pt}
    \IfFileExists{../results/figures/#1_time.pdf}
      {\includegraphics[width=0.8\linewidth]{../results/figures/#1_time.pdf}}
      {\missingfigbox{#1_time}}
    \caption{#2}
    \label{#3}
  \end{figure}%
}

% Slides: figures are 6 by 4 inches, so a full-width copy is too tall for a frame
% that also carries text. Height is capped as a fraction of the text height and
% the aspect ratio kept; crowded frames pass a smaller fraction.
\newcommand{\resultgraphic}[2][0.72]{%
  \IfFileExists{../results/figures/slides/#2.pdf}
    {\includegraphics[width=\linewidth,height=#1\textheight,keepaspectratio]{../results/figures/slides/#2.pdf}}
    {\IfFileExists{../results/figures/#2.pdf}
      {\includegraphics[width=\linewidth,height=#1\textheight,keepaspectratio]{../results/figures/#2.pdf}}
      {\missingfigbox{#2}}}%
}


\usepackage{titlesec}
\usepackage{tocloft}
\usepackage[hidelinks]{hyperref}
\usepackage{algorithm}
\usepackage{algpseudocode}
% polyglossia translates figure and table names for Vietnamese but not the ones
% the algorithm package adds, so the float name and the list heading are set here.
\floatname{algorithm}{Thuật toán}
\renewcommand{\listalgorithmname}{Danh sách thuật toán}
\usepackage[backend=biber, style=numeric, sorting=none]{biblatex}
\addbibresource{refs.bib}

% ---------------------------------------------------------------------------
% Headings. A display block rather than the one-line Roman numeral of the
% Discrete Mathematics template: the word "Chương" and the number sit above a
% rule in small sans, and the title follows underneath. The number stays in
% Arabic so that cross-references read "hình 7.2" without a second numbering
% scheme to keep in step.
% ---------------------------------------------------------------------------
\titleformat{\chapter}[display]
  {\normalfont\sffamily\color{ink}}
  {\color{accent}\sffamily\bfseries\large Chương \thechapter}
  {0.4em}
  {\Huge\bfseries\raggedright}
  [\vspace{0.5em}{\color{accent}\titlerule[1.2pt]}]
\titlespacing*{\chapter}{0pt}{0pt}{28pt}

\titleformat{\section}{\normalfont\sffamily\Large\bfseries\color{ink}}{\thesection}{0.7em}{}
\titleformat{\subsection}{\normalfont\sffamily\large\bfseries\color{ink}}{\thesubsection}{0.6em}{}

\renewcommand{\cftchapfont}{\sffamily\bfseries}
\renewcommand{\cftchappagefont}{\sffamily\bfseries}
\renewcommand{\cftsecfont}{\sffamily}
\setlength{\cftchapnumwidth}{2.2em}

\theoremstyle{definition}
\newtheorem{dinhly}{Định lý}[chapter]
\newtheorem{menhde}[dinhly]{Mệnh đề}
\newtheorem{nhanxet}[dinhly]{Nhận xét}

\begin{document}

\begin{titlepage}
  \centering
  \vspace*{2cm}
  {\sffamily\large\color{accent} TỐI ƯU HÓA NÂNG CAO\par}
  \vspace{0.6cm}
  {\color{accent}\rule{0.55\linewidth}{1.2pt}\par}
  \vspace{1.4cm}
  {\sffamily\Huge\bfseries Tối ưu hóa hàm mất mát Ridge\par}
  \vspace{0.5cm}
  {\sffamily\Large cho bài toán dự đoán lãi suất khoản vay\par}
  \vspace{0.8cm}
  {\large Bộ dữ liệu Lending Club, 2007--2018\par}
  \vfill
  {\large Nhóm \dots\par}
  {\large Lớp Khoa học dữ liệu\par}
  \vspace{1cm}
  {\large \today\par}
\end{titlepage}

\chapter*{Tóm tắt}
\addcontentsline{toc}{chapter}{Tóm tắt}

Báo cáo so sánh bốn thuật toán tối ưu hóa tự cài đặt, gradient descent đầy đủ,
mini-batch SGD, gradient tăng tốc Nesterov và phương pháp Newton, mỗi thuật toán
với cả bước cố định lẫn backtracking, trên cùng một hàm mất mát Ridge dựng từ
\num{1200000} khoản vay Lending Club với $d = 116$ thuộc tính. Hàm mục tiêu lồi
mạnh và có nghiệm đóng \cite{boyd2004convex}, nên giá trị tối ưu $\fstar$ tính được tới sai số máy và
mọi biểu đồ hội tụ vẽ sai số $f(w_k) - \fstar$ trên thang logarit.

Ở cấu hình tốt nhất của từng phương pháp, Newton đạt ngưỡng $10^{-6}$ sau một
vòng lặp và khoảng $0{,}25$ giây, gradient tăng tốc cần 75 vòng và $1{,}05$ giây,
backtracking cần 143 vòng và $3{,}14$ giây, còn gradient descent với bước cố định
tốt nhất cần 331 vòng và $4{,}75$ giây. SGD nhanh nhất trên mỗi đơn vị thời gian
nhưng dừng ở sai số $6{,}0 \cdot 10^{-4}$ vì bước hằng chỉ đưa dãy lặp vào một
lân cận của nghiệm.

Ba quan sát đi ngược dự đoán ban đầu và được giải thích trong báo cáo: bước
$1{,}9/L$ nhanh gấp hơn hai lần bước tối ưu $2/(L+\mu)$ của lý thuyết, do sai số
ban đầu dồn vào các hướng trị riêng lớn; backtracking thắng bước cố định trên cả
hai trục đo chứ không chỉ trên trục vòng lặp; và tốc độ giảm của bước SGD quyết
định kết luận về quy tắc chọn bước, tới mức hai giá trị hợp lý cho hai kết luận
ngược nhau.

\tableofcontents

% ===========================================================================
\chapter{Bài toán và ba hằng số}
% ===========================================================================

\section{Hàm mục tiêu}

Bài toán hồi quy Ridge \cite{hoerl1970ridge} trên $n$ điểm dữ liệu và $d$ thuộc
tính là bài toán cực tiểu hóa không ràng buộc
\begin{equation}\label{eq:objective}
  \min_{w \in \R^{d}} \quad f(w) = \frac{1}{2n} \norm{Xw - y}^{2} + \frac{\lambda}{2} \norm{w}^{2},
\end{equation}
với $X \in \R^{n \times d}$ là ma trận thiết kế đã chuẩn hóa, $y \in \R^{n}$ là
vector lãi suất đã trừ trung bình, và $\lambda > 0$ là hệ số hiệu chỉnh. Hệ số
chặn không xuất hiện trong \eqref{eq:objective} vì đã bị khử bằng cách trừ trung
bình khỏi cả $X$ và $y$; giữ nó lại mà không phạt cũng cho cùng nghiệm, nhưng làm
công thức gradient và Hessian dài thêm mà không đổi nội dung.

Hệ số $\frac{1}{2n}$ thay cho $\frac{1}{2}$ có một lý do cụ thể cho phần so sánh
sau này. Không có $\frac{1}{n}$ thì hằng số Lipschitz tăng tuyến tính theo số
mẫu, nên độ dài bước dò được trên mẫu \num{200000} điểm sẽ không dùng lại được
trên mẫu \num{1200000} điểm, và mọi cấu hình phải quét lại từ đầu ở quy mô lớn.

Gradient và Hessian của \eqref{eq:objective} là
\begin{equation}\label{eq:grad}
  \nabla f(w) = \frac{1}{n} X\trans (Xw - y) + \lambda w,
  \qquad
  \nabla^{2} f(w) = \frac{1}{n} X\trans X + \lambda I.
\end{equation}
Hessian trong \eqref{eq:grad} không chứa $w$, vì $f$ là hàm bậc hai. Tính chất
này chi phối toàn bộ chương~\ref{ch:newton}, vì khai triển Taylor bậc hai của $f$
không còn số hạng dư và bước Newton xuất phát từ $w_{0} = 0$ do đó trùng đúng với
nghiệm đóng
\begin{equation}\label{eq:closed}
  \wstar = \left( \frac{1}{n} X\trans X + \lambda I \right)^{-1} \frac{1}{n} X\trans y .
\end{equation}

\section{Ba hằng số chi phối tốc độ hội tụ}

Ký hiệu $\lambda_{\max}$ và $\lambda_{\min}$ là trị riêng lớn nhất và nhỏ nhất
của ma trận Gram $\frac{1}{n} X\trans X$. Hằng số Lipschitz của gradient, hệ số
lồi mạnh và số điều kiện là
\begin{equation}\label{eq:constants}
  L = \lambda_{\max} + \lambda, \qquad \mu = \lambda_{\min} + \lambda, \qquad \kappa = \frac{L}{\mu}.
\end{equation}
Ba đại lượng này đo được chứ không phải giả định, vì $d = 116$ đủ nhỏ để tính
trọn phổ của ma trận Gram trong dưới một giây. Bảng~\ref{tab:hangso} ghi giá trị
ở hai quy mô mẫu.

\begin{table}[tbp]
  \centering
  \caption{Ba hằng số của bài toán ở hai quy mô mẫu, với $\lambda = 0{,}031623$
    và $d = 116$ cho cả hai.}
  \label{tab:hangso}
  \begin{tabular}{lrrrrr}
    \toprule
    Quy mô & $n$ & $L$ & $\mu$ & $\kappa$ & $\fstar$ \\
    \midrule
    Quét tham số & \num{200000}  & 9{,}1156 & 0{,}034073 & 267{,}5 & 6{,}151114 \\
    Toàn phần    & \num{1200000} & 9{,}0620 & 0{,}034036 & 266{,}2 & 6{,}142806 \\
    \bottomrule
  \end{tabular}
\end{table}

Ba hằng số ở hai quy mô lệch nhau dưới $0{,}6$ phần trăm, cụ thể $0{,}59$ phần
trăm với $L$, $0{,}11$ phần trăm với $\mu$ và $0{,}48$ phần trăm với $\kappa$.
Hệ số $\frac{1}{n}$ trong \eqref{eq:objective} khử ảnh hưởng của kích thước mẫu
lên phổ, nên hai bài toán tuy khác nhau về số điểm lại có cùng độ khó về mặt tối
ưu hóa. Nhờ vậy toàn bộ phần quét lưới tham số chạy được
trên mẫu nhỏ, và chỉ cấu hình tốt nhất mới phải chạy lại ở quy mô đầy đủ.

Nghiệm đóng \eqref{eq:closed} giải bằng phân rã Cholesky cho
$\norm{\nabla f(\wstar)} = 7{,}4 \cdot 10^{-13}$, tức ở mức sai số máy. Nhờ con số
đó, báo cáo đo được sai số tuyệt đối $f(w_k) - \fstar$ của từng thuật toán thay vì
chỉ so sánh các đường cong với nhau. Nếu đổi sang hàm mất mát không có nghiệm đóng, chẳng hạn Huber hay Lasso,
thì $\fstar$ phải ước lượng bằng chính lần chạy dài nhất, và sai số của mọi
sai số của đường tốt nhất chặn dưới mọi đường còn lại.

\section{Cách đo sai số}

Vẽ $f(w_k) - \fstar$ bằng phép trừ trực tiếp có một giới hạn số học đáng kể. Hai
số gần bằng nhau khi trừ mất hết chữ số có nghĩa ở phần thấp, nên với
$\fstar \approx 6{,}15$ thì mọi giá trị dưới khoảng $10^{-16} \fstar$ đều là
nhiễu làm tròn. Vì $f$ là hàm bậc hai và $\nabla f(\wstar) = 0$, khai triển
Taylor không còn số hạng dư và đẳng thức
\begin{equation}\label{eq:gap}
  f(w) - \fstar = \tfrac{1}{2} (w - \wstar)\trans \nabla^{2} f \, (w - \wstar)
\end{equation}
đúng chính xác chứ không phải xấp xỉ. Báo cáo dùng vế phải của \eqref{eq:gap}.

Trên bài toán thử với $n = 400$ và $d = 12$, tại
$\norm{w - \wstar} \approx 10^{-8}$ phép trừ trực tiếp chỉ còn đúng ba chữ số, và tại $10^{-12}$ nó trả về $-1{,}4 \cdot 10^{-17}$, tức một giá
trị âm không có ý nghĩa, trong khi \eqref{eq:gap} cho $3{,}7 \cdot 10^{-24}$.
Trục tung của mọi hình hội tụ trong báo cáo do đó xuống được sâu hơn tám bậc so
với cách vẽ thông thường. Cách này chỉ dùng được vì $f$ bậc hai; với hàm tổng
quát thì \eqref{eq:gap} chỉ còn là xấp xỉ bậc hai quanh $\wstar$.

% ===========================================================================
\chapter{Dữ liệu: rò rỉ và chuẩn hóa}
% ===========================================================================

\section{Bộ dữ liệu và các cột phải loại}

Bộ dữ liệu Lending Club gồm \num{2260668} khoản vay được duyệt trong giai đoạn
2007--2018 với 151 cột, sau khi loại 33 dòng thiếu biến mục tiêu. Biến cần dự
đoán là lãi suất \texttt{int\_rate}, nằm trong khoảng từ $5{,}31$ tới $30{,}99$
phần trăm với độ lệch chuẩn $4{,}83$.

Cột \texttt{sub\_grade} phải loại trước mọi bước khác. Hồi quy \texttt{int\_rate}
theo riêng \texttt{sub\_grade} cho $R^{2} = 0{,}9554$ trên toàn bộ dữ liệu, và
$R^{2} = 0{,}9834$ khi xét riêng các khoản vay phát hành năm 2016, còn
\texttt{grade} cho $R^{2} = 0{,}9088$. Lending Club ấn định lãi suất từ một bảng
tra theo bậc tín dụng, nên hai cột này chính là biến mục tiêu mang tên khác; con số trong phạm vi một năm cao hơn vì bảng tra ấy cố định trong
năm và thay đổi giữa các năm. Cột \texttt{installment} cũng bị loại, vì nó tính
ra từ \texttt{loan\_amnt}, kỳ hạn và chính \texttt{int\_rate} bằng công thức trả
góp. Sau khi loại cả nhóm cột ghi nhận sau giải ngân, mô hình đạt
$R^{2} \approx 0{,}486$ trên tập kiểm tra. Chênh lệch giữa $0{,}955$ và $0{,}486$
là thước đo trực tiếp của phần rò rỉ.

Với riêng phần tối ưu hóa, rò rỉ không làm sai một phép tính nào, vì hàm mục tiêu
vẫn lồi mạnh và mọi thuật toán vẫn hội tụ về cùng một $\wstar$. Nó chỉ làm hỏng
phần đối chiếu chất lượng dự báo ở chương~\ref{ch:lambda}, nơi RMSE trên tập
kiểm tra được dùng để cân nhắc hệ số hiệu chỉnh.

\section{Một hướng suy biến do cách ghi dữ liệu}

Lần dựng ma trận thiết kế đầu tiên cho $d = 117$ với trị riêng nhỏ nhất của ma
trận Gram bằng $1{,}07 \cdot 10^{-7}$, kéo theo $\mu = \lambda$ đúng tới sáu chữ
số. Vector riêng ứng với trị riêng đó có đúng hai hệ số khác không, bằng
$+0{,}7071$ tại \texttt{fico\_range\_low} và $-0{,}7071$ tại
\texttt{fico\_range\_high}, tức đúng hiệu hai cột chia cho $\sqrt{2}$.

Truy lại dữ liệu gốc thì \texttt{fico\_range\_high} bằng
\texttt{fico\_range\_low} cộng $4$ ở \num{2260227} trên \num{2260668} dòng, với
hệ số tương quan $0{,}99999991$. Hai cột lệch nhau đúng một hằng số nên trùng
nhau hoàn toàn sau khi chuẩn hóa, và ma trận Gram nhận một hướng gần suy biến
không mang thông tin nào. Đã loại \texttt{fico\_range\_high} cùng nhóm với
\texttt{funded\_amnt} và \texttt{funded\_amnt\_inv}, và $d$ giảm còn 116. Sau khi
loại, trị riêng nhỏ nhất là $2{,}45 \cdot 10^{-3}$, nên $\mu = 0{,}0341$ vẫn do
$\lambda = 0{,}0316$ chi phối tới 93 phần trăm nhưng không còn suy biến.
Hình~\ref{fig:spectrum} vẽ phổ trị riêng cùng với sàn do $\lambda$ đặt.

\resultfig{spectrum}{Phổ trị riêng của ma trận Gram $\frac{1}{n}X\trans X$ và sàn
  do hệ số hiệu chỉnh đặt. Mọi trị riêng nằm dưới đường ngang đều bị $\lambda$
  thay thế trong công thức $\mu = \lambda_{\min} + \lambda$.}{fig:spectrum}

\section{Chuẩn hóa cột và số điều kiện}

Bảng~\ref{tab:scaling} so sánh ba cách mã hóa cùng một tập dòng: giữ nguyên thang
gốc, chỉ trừ trung bình, và chuẩn hóa đầy đủ về trung bình $0$ độ lệch chuẩn $1$.

\begin{table}[tbp]
  \centering
  \caption{Ảnh hưởng của cách xử lý cột lên số điều kiện, đo với
    $\lambda = 0{,}031623$ và $n = \num{200000}$. Cột cuối là số vòng lặp
    gradient descent cần để đạt $f - \fstar < 10^{-6}$ với bước $1/L$.}
  \label{tab:scaling}
  \begin{tabular}{lrrrr}
    \toprule
    Cách mã hóa & $L$ & $\mu$ & $\kappa$ & Số vòng lặp \\
    \midrule
    Thô                & $1{,}225 \cdot 10^{11}$ & 0{,}031628 & $3{,}87 \cdot 10^{12}$ & không đạt \\
    Chỉ trừ trung bình & $6{,}039 \cdot 10^{10}$ & 0{,}031628 & $1{,}91 \cdot 10^{12}$ & không đạt \\
    Chuẩn hóa đầy đủ   & 9{,}1156                & 0{,}034073 & 267{,}5                & 630 \\
    \bottomrule
  \end{tabular}
\end{table}

Số điều kiện giảm $1{,}45 \cdot 10^{10}$ lần khi chuyển từ thang gốc sang thang
chuẩn hóa, nhưng phần đóng góp của hai thao tác rất chênh lệch. Trừ trung bình
chỉ hạ $\kappa$ đúng hai lần, còn chia độ lệch chuẩn hạ tiếp $7 \cdot 10^{9}$
lần. Các cột đứng ở những thang khác nhau tới sáu bậc, \texttt{tot\_cur\_bal} cỡ
$10^{6}$ bên cạnh \texttt{dti} cỡ hàng chục, nên cột có đơn vị lớn nhất chi phối
trị riêng lớn nhất; trừ đi trung bình không đụng
gì tới chênh lệch thang đo đó.

Ở hai biến thể chưa chia thang, $\mu$ bằng đúng $\lambda$ tới năm chữ số, tức ma
trận Gram suy biến về mặt số học khi đặt cạnh $L$ cỡ $10^{11}$. Hình~\ref{fig:scaling} cho thấy hệ quả: sau \num{4000} vòng lặp, gradient descent
mới hạ được sai số từ $5{,}53$ xuống $3{,}74$, tức chưa tới một phần ba. Chuẩn hóa vì thế
không phải một thói quen tiền xử lý mà là một can thiệp trực tiếp lên độ khó của
bài toán tối ưu hóa. Kết luận này mất hiệu lực nếu mọi cột vốn đã cùng thang, khi
đó phép chia không còn gì để sửa.

\resultpair{scaling-kappa}{Gradient descent với bước $1/L$ trên ba cách mã hóa
  cột. Hai đường trên cùng gần như nằm ngang trong suốt \num{4000} vòng
  lặp.}{fig:scaling}

\section{Chọn hệ số hiệu chỉnh}

Hệ số $\lambda$ chọn bằng cross-validation $5$ fold trên tập huấn luyện, quét trên
lưới logarit từ $10^{-6}$ tới $10^{2}$. Đường cong sai số phẳng trên nhiều bậc
$\lambda$, nên lấy điểm cực tiểu là lựa chọn tùy tiện. Cực tiểu của một đường phẳng rơi vào
chỗ nhiễu đặt nó, và ở đây nó rơi vào $\lambda = 0{,}001778$, tức
gần đầu nhỏ nhất của lưới và do đó ứng với $\kappa$ lớn nhất.

Báo cáo dùng quy tắc một sai số chuẩn \cite{hastie2009elements}, tức chọn
$\lambda$ lớn nhất còn nằm trong một sai số chuẩn của giá trị tốt nhất. Kết quả
là $\lambda = 0{,}031623$, lớn hơn 18 lần so với điểm cực tiểu, đổi lại sai số
cross-validation tăng từ $12{,}0430$ lên $12{,}0607$, tức $0{,}15$ phần trăm.
Chương~\ref{ch:lambda} cho thấy khoản đánh đổi này mua được gì về phía tối ưu
hóa. Quy tắc sẽ không phù hợp nếu đường cong cross-validation có cực tiểu nhọn,
khi đó dịch $\lambda$ ra xa cực tiểu phải trả giá thật về chất lượng dự báo.

% ===========================================================================
\chapter{Cài đặt và kiểm thử}
% ===========================================================================

\section{Hai trục của bài toán}

Đề bài yêu cầu bốn thuật toán, mỗi thuật toán hai cơ chế chọn bước, tức một ma
trận bốn nhân hai. Mã nguồn tách đúng hai trục đó thành hai loại đối tượng rời
nhau. Trục thứ nhất trả lời câu hỏi đi từ điểm nào theo hướng nào, gồm bốn lớp
\texttt{SteepestDescent}, \texttt{MiniBatch}, \texttt{Nesterov} và
\texttt{NewtonStep}. Trục thứ hai trả lời câu hỏi đi bao xa, gồm ba lớp
\texttt{Fixed}, \texttt{Armijo} và \texttt{Decay}. Một hàm lặp duy nhất ghép hai
trục lại, giữ đồng hồ, ghi log và kiểm tra điều kiện dừng.

Cách tách này có ba hệ quả thực tế. Tám cấu hình bắt buộc trở thành tích Descartes
của hai danh sách, nên câu hỏi đã cài đủ hay chưa trả lời được bằng một vòng lặp
trong bộ kiểm thử thay vì bằng cách đọc lại mã. Quy tắc dừng đồng hồ trước khi
ghi log chỉ phải kiểm tra ở một chỗ thay vì bốn. Và điểm xuất phát của bước được
tách khỏi dãy lặp hiện tại, vì Nesterov tính gradient tại điểm ngoại suy $y_k$
chứ không tại $w_k$; gộp hai khái niệm này là lỗi thường gặp khi cài Nesterov.

Cách tách không phủ hết mọi thuật toán. Adam co giãn theo từng tọa độ nên không
có độ dài bước vô hướng, và nếu bổ sung thì phần co giãn phải nằm trong lớp hướng
đi, làm nhòe ranh giới hai trục.

\section{Điều kiện Armijo ở dạng tổng quát}

Thuật toán~\ref{alg:armijo} là backtracking line search theo điều kiện giảm đủ
của Armijo \cite{armijo1966minimization}.

\begin{algorithm}
  \caption{Backtracking line search (Armijo)}
  \label{alg:armijo}
  \begin{algorithmic}[1]
    \Require descent direction $p$ at $y$, parameters $c \in (0,1)$, $\rho \in (0,1)$, $t_{0} > 0$
    \State $t \gets t_{0}$
    \While{$f(y + t p) > f(y) + c\, t\, \nabla f(y)\trans p$}
      \State $t \gets \rho\, t$
    \EndWhile
    \State \Return $t$
  \end{algorithmic}
\end{algorithm}

Điều kiện dừng viết theo tích vô hướng $\nabla f(y)\trans p$ chứ không theo
$-\norm{\nabla f(y)}^{2}$. Hai dạng trùng nhau khi $p = -\nabla f(y)$, nhưng chỉ
dạng tổng quát còn đúng với hướng Newton, nơi
$p = -(\nabla^{2} f)^{-1} \nabla f$ và do đó
$\nabla f\trans p = -\nabla f\trans (\nabla^{2} f)^{-1} \nabla f < 0$. Vì lớp
\texttt{Armijo} nhận hướng từ bên ngoài và không biết ai tạo ra nó, nó buộc phải
dùng dạng tổng quát.

\section{Bốn phép kiểm tra}

Trước khi chạy lưới tham số, mã nguồn qua bốn phép kiểm tra. So sánh gradient
giải tích với sai phân trung tâm cho sai số tương đối $1{,}2 \cdot 10^{-9}$, và
so sánh Hessian với sai phân của gradient cho $1{,}8 \cdot 10^{-9}$; cả hai đều
dưới ngưỡng $10^{-6}$ mà bậc $\varepsilon^{2}$ của sai phân trung tâm cho phép
kỳ vọng. Nghiệm đóng cho $\norm{\nabla f(\wstar)} = 2{,}7 \cdot 10^{-15}$ trên bài
toán thử. Phép kiểm cuối chạy toàn bộ tám cấu hình trên bài toán nhỏ có nghiệm
biết trước và xác nhận cả tám đều về cùng một $\wstar$ với sai số dưới $10^{-8}$.

Một chi tiết cài đặt cần nêu vì nó từng gây lỗi. Công thức momentum
$\beta = (\sqrt{\kappa}-1)/(\sqrt{\kappa}+1)$ chỉ đúng khi $t = 1/L$, nên khi
ghép với backtracking phải tính lại theo bước thật bằng
$\beta = (1 - \sqrt{t\mu})/(1 + \sqrt{t\mu})$. Lớp \texttt{Nesterov} nhận đúng
giá trị $t$ mà line search vừa chọn qua hàm \texttt{accept}, nên hai đại lượng
không có đường nào để lệch nhau. Vì $t$ chỉ biết sau khi line search chạy xong
còn $\beta$ phải có trước để dựng $y_k$, cài đặt dùng bước của vòng trước và lấy
$1/L$ cho vòng đầu; với bước cố định thì hai giá trị trùng nhau.

% ===========================================================================
\chapter{Chọn độ dài bước thế nào}\label{ch:step}
% ===========================================================================

\section{Bước cố định và ngưỡng phân kỳ}

Lý thuyết cho biết gradient descent với bước hằng hội tụ khi $0 < t < 2/L$, và
với hàm bậc hai lồi mạnh thì bước tối ưu là $t = 2/(L+\mu)$
\cite{nesterov2018lectures}. Bảng~\ref{tab:stepfixed} ghi kết quả quét tám giá
trị quanh ngưỡng đó.

\begin{table}[tbp]
  \centering
  \caption{Gradient descent với bước cố định, $n = \num{200000}$,
    $\kappa = 267{,}5$. Cột giữa là số vòng lặp cần để đạt
    $f - \fstar < 10^{-6}$.}
  \label{tab:stepfixed}
  \begin{tabular}{lrrl}
    \toprule
    Bước & Số vòng lặp & Thời gian (s) & Trạng thái \\
    \midrule
    $t = 2{,}1/L$    & --   & --      & phân kỳ ở vòng 94 \\
    $t = 2/L$        & --   & --      & đình trệ ở $9{,}4 \cdot 10^{-2}$ \\
    $t = 1{,}9/L$    & 331  & 4{,}75  & hội tụ \\
    $t = 2/(L+\mu)$  & 766  & 10{,}87 & hội tụ \\
    $t = 1/L$        & 630  & 8{,}98  & hội tụ \\
    $t = 0{,}5/L$    & 1261 & 18{,}04 & chạm giới hạn vòng lặp \\
    $t = 0{,}1/L$    & --   & --      & không đạt sau \num{5000} vòng \\
    \bottomrule
  \end{tabular}
\end{table}

Ngưỡng $2/L$ hiện ra đúng như lý thuyết mô tả. Với $t = 2{,}1/L$, hệ số co dọc
hướng dốc nhất là $\abs{1 - tL} = 1{,}1 > 1$, nên thành phần sai số theo hướng đó
nở ra mỗi vòng và dãy lặp phân kỳ sau 94 vòng. Với $t = 2/L$ đúng bằng ngưỡng, hệ
số bằng $1$ nên thành phần đó dao động mà không tắt, và sai số dừng ở
$9{,}4 \cdot 10^{-2}$. Hình~\ref{fig:stepfixed} vẽ cả ba trường hợp.

\resultpair{step-fixed}{Gradient descent với các bước cố định khác nhau. Đường
  phân kỳ ứng với $t = 2{,}1/L$ được giữ lại trong hình chứ không loại
  bỏ.}{fig:stepfixed}

\section{Chênh lệch giữa \texorpdfstring{$1{,}9/L$}{1,9/L} và bước tối ưu lý thuyết}

Bước $2/(L+\mu)$ cần 766 vòng lặp trong khi bước $1{,}9/L$ chỉ cần 331, tức chậm
hơn $2{,}3$ lần dù lý thuyết gọi nó là bước tối ưu. Lý thuyết không sai; nó trả lời một câu hỏi khác.

Cận $\left((\kappa-1)/(\kappa+1)\right)^{2k}$ là cận cực tiểu hóa trường hợp xấu
nhất trên mọi điểm khởi tạo, nên bước tối ưu cân bằng hai đầu phổ, tức là với
$t = 2/(L+\mu)$ thì $\abs{1 - tL} = 0{,}9926$ và $\abs{1 - t\mu} = 0{,}99255$,
gần bằng nhau. Bước $1{,}9/L$ hy sinh một chút ở đầu nhỏ, $\abs{1-t\mu} =
0{,}99290$, để đổi lấy $\abs{1 - tL} = 0{,}9000$ ở đầu lớn.

Điểm khởi tạo $w_{0} = 0$ không phải trường hợp xấu nhất. Phân tích sai số ban
đầu theo cơ sở riêng của Hessian cho thấy các hướng có trị riêng lớn hơn $1$ mang
$80{,}5$ phần trăm của $f(w_{0}) - \fstar$, trong khi các hướng có trị riêng nhỏ
hơn $0{,}1$ chỉ mang $0{,}3$ phần trăm, và riêng hướng ứng với trị riêng $3{,}76$
chiếm $35{,}9$ phần trăm. Chính dạng toàn phương \eqref{eq:gap} gây ra sự chênh
lệch đó. Viết trong cơ sở riêng thì
\begin{equation}\label{eq:modal}
  f(w_k) - \fstar = \frac{1}{2} \sum_{i=1}^{d} \lambda_{i} c_{i}^{2} \left( 1 - t \lambda_{i} \right)^{2k},
\end{equation}
nên mỗi hướng được cân theo đúng trị riêng của nó, và các hướng dốc đóng góp
nhiều gấp bội vào sai số cần triệt tiêu.

Mô hình \eqref{eq:modal} dự đoán được con số chứ không chỉ chiều hướng: nó cho
sai số $9{,}80 \cdot 10^{-7}$ tại vòng 331 với bước $1{,}9/L$, và
$9{,}95 \cdot 10^{-7}$ tại vòng 766 với bước $2/(L+\mu)$, khớp với ngưỡng
$10^{-6}$ ở đúng hai vòng lặp quan sát được. Kết luận đảo chiều khi cần độ chính
xác rất cao: lúc đó mọi hướng dốc đã tắt và tốc độ do hướng ứng với $\mu$ quyết
định, nơi $2/(L+\mu)$ nhỉnh hơn. Với ngưỡng $10^{-6}$ mà báo cáo dùng, giai đoạn
đó chưa tới.

\section{Dò bước mà không tính \texorpdfstring{$L$}{L}}

Đề bài gợi ý thử các bước $10^{-3}$, $10^{-2}$, $10^{-1}$, $1$ và $10$. Trong lưới
đó chỉ $t = 0{,}1$ hội tụ, sau 691 vòng lặp; $t = 0{,}01$ không đạt ngưỡng sau
\num{5000} vòng, còn $t = 1$ phân kỳ sau 5 vòng và $t = 10$ sau 2 vòng. Vùng hội
tụ là khoảng $(0;\ 0{,}219)$, nên lưới cách nhau một bậc chỉ có đúng một điểm rơi
vào đó. Tính $L$ và $\mu$ tốn dưới một giây, còn quét năm giá trị tốn 195 giây và
vẫn không tìm được bước gần tối ưu.

\section{Backtracking}

Bảng~\ref{tab:armijo} ghi sáu cấu hình backtracking tốt nhất trong lưới 18 cấu
hình đã quét.

\begin{table}[tbp]
  \centering
  \caption{Backtracking line search, sáu cấu hình tốt nhất trong 18 cấu hình đã
    quét. Cột cuối là số lần đánh giá hàm mục tiêu trung bình mỗi vòng lặp.}
  \label{tab:armijo}
  \begin{tabular}{llrrr}
    \toprule
    $c$ & $\rho$, $t_{0}$ & Số vòng lặp & Thời gian (s) & Lần đánh giá $f$ \\
    \midrule
    0{,}3    & $\rho = 0{,}5$, $t_{0} = 1$    & 143 & 3{,}14 & 4{,}94 \\
    0{,}1    & $\rho = 0{,}5$, $t_{0} = 1$    & 149 & 3{,}23 & 4{,}04 \\
    $10^{-4}$ & $\rho = 0{,}5$, $t_{0} = 1$    & 167 & 3{,}27 & 3{,}95 \\
    0{,}3    & $\rho = 0{,}5$, $t_{0} = 10/L$ & 190 & 4{,}34 & 4{,}17 \\
    0{,}3    & $\rho = 0{,}8$, $t_{0} = 1$    & 312 & 9{,}62 & 7{,}77 \\
    0{,}3    & $\rho = 0{,}9$, $t_{0} = 1$    & \multicolumn{1}{c}{--} & \multicolumn{1}{c}{--} & 12{,}4 \\
    \bottomrule
  \end{tabular}
\end{table}

Backtracking thắng bước cố định trên cả hai trục đo, chứ không phải chỉ thắng về
số vòng lặp và thua về thời gian như thường được mô tả: 143 vòng và $3{,}14$ giây
so với 331 vòng và $4{,}75$ giây của bước cố định tốt nhất, dù mỗi vòng tốn thêm
$4{,}94$ lần đánh giá hàm mục tiêu. Cơ chế nằm ở chỗ điều kiện Armijo không so
với $L$ toàn cục mà so với thương Rayleigh tại chỗ: bước được chấp nhận thỏa
$t \le 2(1-c) \norm{g}^{2} / (g\trans \nabla^{2} f\, g)$, nên khi gradient nghiêng
về các hướng phẳng thì mẫu số nhỏ và line search nhận bước lớn hơn $2/L$ rất
nhiều. Bước cố định không làm được điều đó vì nó phải an toàn cho hướng dốc nhất
ở mọi vòng lặp. Kết luận đảo chiều khi mỗi lần đánh giá $f$ đắt ngang một lần
tính gradient và $\kappa$ nhỏ, khi đó phần chi phí thêm không còn được bù lại.

Hai tham số của line search không tương đương nhau về mức ảnh hưởng. Toàn bộ sáu
cấu hình $\rho = 0{,}5$ đứng trên toàn bộ mười hai cấu hình $\rho = 0{,}8$ và
$\rho = 0{,}9$, trong khi $c$ đổi từ $10^{-4}$ lên $0{,}3$, tức 3000 lần, chỉ làm
số vòng lặp nhích từ 167 xuống 143. Tham số $\rho$ quyết định độ thô của lưới
bước thử. Đi từ $t_{0}$ xuống cùng một bước, $\rho = 0{,}9$ cần gấp khoảng bảy
lần số lần thử so với $\rho = 0{,}5$, và mỗi lần thử tốn một lần tính $f$.

Hình~\ref{fig:armijo} vẽ các cấu hình này theo hai trục đo, còn
hình~\ref{fig:armijocost} vẽ số lần đánh giá hàm theo vòng lặp và cho thấy chi
phí không phân bố đều. Ở giai đoạn đầu, line search nhận bước ngay ở lần thử thứ
nhất hoặc thứ hai. Về cuối, khi $f - \fstar$ chạm sàn phân giải số học của $f$,
mức giảm thật nhỏ hơn nhiễu làm tròn nên mọi bước thử đều trượt điều kiện, và số
lần đánh giá vọt lên hàng chục. Phần đuôi đó là chi phí chết, và tiêu chí dừng
theo đình trệ cắt được nó.

\resultfig{step-armijo_cost}{Số lần đánh giá hàm mục tiêu mỗi vòng lặp, cho các
  cấu hình backtracking với $t_{0} = 1$.}{fig:armijocost}

\resultpair{step-armijo}{Backtracking với các cấu hình $c$ và $\rho$ khác
  nhau.}{fig:armijo}

% ===========================================================================
\chapter{Quán tính: đổi bộ nhớ lấy số vòng lặp}\label{ch:momentum}
% ===========================================================================

Ba chương tiếp theo đặt bốn phương pháp vào cùng một khung đánh đổi giữa số vòng
lặp và giá của mỗi vòng. Chương này xét cách rẻ nhất: giữ thêm một vector trong
bộ nhớ mà không tăng chi phí mỗi vòng.

\section{Ba công thức momentum}

Gradient tăng tốc Nesterov cập nhật theo
\begin{equation}\label{eq:agd}
  y_k = w_k + \beta_k (w_k - w_{k-1}), \qquad w_{k+1} = y_k - t\, \nabla f(y_k),
\end{equation}
Cập nhật \eqref{eq:agd} vẫn dùng đúng một lần tính gradient mỗi vòng như gradient descent, chỉ thêm một
vector $w_{k-1}$ trong bộ nhớ. Bảng~\ref{tab:momentum} so sánh ba cách chọn
$\beta_k$ với cùng bước $t = 1/L$.

\begin{table}[tbp]
  \centering
  \caption{Ba công thức momentum, cùng bước $t = 1/L$, so với gradient descent.}
  \label{tab:momentum}
  \begin{tabular}{lrrl}
    \toprule
    Công thức & Số vòng lặp & Thời gian (s) & Trạng thái \\
    \midrule
    Gradient descent                        & 3000 & 42{,}13 & chưa hội tụ \\
    $\beta$ theo $t$ và $\mu$               & 286  & 4{,}01  & hội tụ \\
    $\beta = 0{,}9$                         & 314  & 4{,}42  & hội tụ \\
    $\beta_k = (k-1)/(k+2)$                 & 894  & 12{,}78 & đình trệ ở $2{,}1 \cdot 10^{-11}$ \\
    \bottomrule
  \end{tabular}
\end{table}

Quán tính đổi một vector bộ nhớ lấy hơn mười lần số vòng lặp: 286 vòng và
$4{,}01$ giây so với 3000 vòng và $42{,}13$ giây của gradient descent trên cùng
bước. Vì giá mỗi vòng của hai phương pháp như nhau, tỷ số thời gian bằng đúng tỷ
số số vòng lặp. Khoảng cách này thu hẹp khi $\kappa$ nhỏ, vì $\sqrt{\kappa}$ và
$\kappa$ tiến lại gần nhau; chương~\ref{ch:lambda} đo cụ thể.

Hằng số $\beta = 0{,}9$ mà đề bài gợi ý cần 314 vòng, chỉ chậm hơn 10 phần trăm
so với công thức lý thuyết. Không nên đọc con số đó thành một kết luận chung: với $\kappa = 267{,}5$ thì $(\sqrt{\kappa}-1)/(\sqrt{\kappa}+1) = 0{,}885$,
tình cờ sát $0{,}9$, và công thức momentum vốn không nhạy quanh cực trị. Với bài
toán có $\kappa$ lệch xa, chẳng hạn $\kappa = 10$ ứng với $\beta^{*} = 0{,}52$,
hằng số $0{,}9$ sẽ sai rõ.

Công thức $\beta_k = (k-1)/(k+2)$ đình trệ ở sai số $2{,}1 \cdot 10^{-11}$ sau
894 vòng. Dãy này tiến tới $1$ nên momentum tích lũy vượt quá mức cần thiết và
sinh dao động tuần hoàn thấy rõ trên hình~\ref{fig:momentum}. Công thức được
thiết kế cho trường hợp không giả định biết $\mu$, nên trên hàm lồi mạnh đã biết
$\mu$ thì nó không có lý do để dùng.

\resultpair{momentum-formula}{Ba công thức momentum so với gradient descent, cùng
  bước $t = 1/L$.}{fig:momentum}

\section{Khởi động lại thích nghi}

Kỹ thuật khởi động lại thích nghi \cite{odonoghue2015adaptive} đặt bộ đếm
momentum về $0$ khi phát hiện
$\nabla f(y_k)\trans (w_{k+1} - w_k) > 0$, tức khi quán tính đang đẩy ngược chiều
gradient. Ba công thức ở bảng~\ref{tab:momentum} phản ứng rất khác nhau. Với $\beta_k$ tăng
dần, số vòng lặp để đạt $10^{-6}$ giảm từ 143
xuống 89 và trạng thái chuyển từ đình trệ sang hội tụ, còn với $\beta$ tính theo
$\mu$ và với $\beta = 0{,}9$ thì kết quả có và không khởi động lại trùng nhau tới
từng chữ số.

Hai kết quả trùng nhau vì điều kiện khởi động lại không kích hoạt lần nào khi
$\beta$ đã đặt đúng theo $\mu$: momentum không bao giờ vượt đà nên không có gì để
sửa. Khởi động lại vì thế là cơ chế bù cho việc không biết $\mu$, chứ không phải
một cải tiến chung. Khi $\mu$ không ước lượng được, nó trở
thành bắt buộc. Hình~\ref{fig:restart} vẽ cả sáu lần chạy.

\resultpair{momentum-restart}{Ba công thức momentum, có và không có khởi động lại
  thích nghi.}{fig:restart}

% ===========================================================================
\chapter{Ngẫu nhiên hóa: đổi độ chính xác lấy giá mỗi vòng}\label{ch:sgd}
% ===========================================================================

\section{Kích thước lô}

Mini-batch SGD \cite{robbins1951stochastic, bottou2018optimization} thay gradient
đầy đủ bằng ước lượng trên một lô ngẫu nhiên $\Batch$ với $\abs{\Batch} = B$:
\begin{equation}\label{eq:sgd}
  g_k = \frac{1}{B} \sum_{i \in \Batch} \left( x_i\trans w_k - y_i \right) x_i + \lambda w_k,
  \qquad w_{k+1} = w_k - \eta_k g_k .
\end{equation}
Chi phí mỗi vòng của \eqref{eq:sgd} giảm $n/B$ lần, nên trục số vòng lặp không so
sánh công bằng được SGD với các phương pháp đầy đủ; hình~\ref{fig:batchsize} vì
thế vẽ thêm theo số epoch.

Hằng số Lipschitz của mất mát một mẫu lớn hơn $L$ rất nhiều, vì $L$ là trung bình
còn nó là cực đại. Độ dài bước an toàn cho lô kích thước $B$ do đó lấy theo
\begin{equation}\label{eq:LB}
  L_{B} = L + \frac{L_{\max} - L}{B}, \qquad L_{\max} = \max_{i} \norm{x_i}^{2} + \lambda .
\end{equation}

Với bước $1/L_{B}$ riêng cho từng $B$, mọi kích thước lô đều dừng quanh sai số
$5 \cdot 10^{-4}$ sau 40 epoch, nhưng chi phí rất khác nhau: $B = 32$ tốn
\num{250000} vòng lặp và $4{,}16$ giây, còn $B = 2048$ tốn 3880 vòng và $2{,}17$
giây. Kích thước lô vì thế đổi giá mỗi vòng chứ không đổi mức sàn, do bước đặt
theo $1/L_{B}$ đã bù trừ giữa phương sai của gradient ước lượng và số bước thực
hiện được.

Dùng chung một bước $1/L$ cho mọi kích thước lô thì cả ba giá trị $B = 32$,
$B = 256$ và $B = 2048$ đều phân kỳ. Đo trên bộ dữ liệu này,
$L_{\max} = \num{200056}$ so với $L = 9{,}12$, tức gấp \num{21947} lần, nên ngay
cả $B = 2048$ cũng chỉ hạ $L_{B}$ xuống $106{,}8$, vẫn lớn hơn $L$ mười lần.

Con số $L_{\max}$ bị đúng một dòng dữ liệu kéo lên. Trung vị của $\norm{x_i}^{2}$
là $82{,}2$ và phân vị 99 là $560{,}9$, trong khi cực đại là \num{200056}; tọa độ
lớn nhất của dòng đó nằm ở cột chỉ dấu \texttt{addr\_state=other} với giá trị 447
độ lệch chuẩn. Chuẩn hóa một cột chỉ dấu có tỷ lệ rất nhỏ biến mỗi giá trị $1$
thành một đỉnh nhọn, nên $L_{\max}$ trong \eqref{eq:LB} ở đây phản ánh cách mã hóa nhiều hơn phản ánh
dữ liệu.

\resultpair{batch-size}{SGD với các kích thước lô khác nhau. Nhóm đường phân kỳ
  ứng với việc dùng chung bước $1/L$ cho mọi $B$.}{fig:batchsize}

\section{Bước khởi đầu}

Dù $L_{B}$ bị ngoại lai kéo lệch, bước $1/L_{B} = 0{,}001265$ tại $B = 256$ vẫn
nằm đúng vùng tốt nhất đo được bằng thực nghiệm. Quét $\eta_{0}$ trên lưới
logarit cho sai số cuối $4{,}21 \cdot 10^{-4}$ tại $\eta_{0} = 10^{-3}$,
$4{,}89 \cdot 10^{-3}$ tại $3 \cdot 10^{-4}$, rồi hỏng hẳn ở $3 \cdot 10^{-3}$
với sai số $1{,}03 \cdot 10^{3}$ và phân kỳ từ $10^{-2}$ trở lên. Cận lý thuyết
do đó vẫn đáng tính: nó là chặn trên an toàn, và ở bài này nó chặt.
Hình~\ref{fig:batcheta} vẽ toàn bộ lưới.

\resultpair{batch-eta}{SGD với các bước khởi đầu khác nhau, $B = 256$.}{fig:batcheta}

\section{Quy tắc giảm bước, và một kết luận phụ thuộc tham số}

SGD với bước hằng không hội tụ về $\wstar$ mà dao động trong một lân cận bán kính
cỡ $\bigO(\eta \sigma^{2} / \mu)$, với $\sigma^{2}$ là phương sai của gradient
ngẫu nhiên \cite{bottou2018optimization}. Trên thang logarit, hiện tượng này hiện
ra thành một đường nằm ngang, và nó là kết quả đúng cần giải thích chứ không phải
lỗi cài đặt. Bảng~\ref{tab:schedule} so sánh bốn quy tắc chọn bước sau 200 epoch.

\begin{table}[tbp]
  \centering
  \caption{Bốn quy tắc chọn bước cho SGD với $B = 256$, sai số sau 200 epoch,
    trung vị của 5 seed. Hai dòng \texttt{inverse} khác nhau ở tốc độ giảm.}
  \label{tab:schedule}
  \begin{tabular}{llr}
    \toprule
    Quy tắc & $\eta_k$ & $f - \fstar$ \\
    \midrule
    Bậc thang, giảm nửa mỗi 40 epoch & $\eta_{0} 2^{-\lfloor k/P \rfloor}$ & $8{,}49 \cdot 10^{-7}$ \\
    Nghịch đảo, $\gamma = \mu \eta_{0}$      & $\eta_{0}/(1+\gamma k)$ & $6{,}94 \cdot 10^{-6}$ \\
    Hằng số                                   & $\eta_{0}$              & $1{,}06 \cdot 10^{-3}$ \\
    Nghịch đảo, $\gamma = 1/\text{epoch}$     & $\eta_{0}/(1+\gamma k)$ & $1{,}46 \cdot 10^{-2}$ \\
    Căn bậc hai                               & $\eta_{0}/\sqrt{k+1}$   & $4{,}58 \cdot 10^{-1}$ \\
    \bottomrule
  \end{tabular}
\end{table}

Quy tắc giảm dần thắng bước hằng khoảng ba bậc, nhưng chỉ khi tốc độ giảm đặt
theo $\mu$. Hai dòng nghịch đảo trong bảng~\ref{tab:schedule} khác nhau đúng một
tham số và cho hai kết luận ngược nhau: với $\gamma = 1/\text{epoch}$, một giá trị
trông tự nhiên, bước tụt xuống $3 \cdot 10^{-6}$ ở epoch 400 và dãy đóng băng ở
sai số $1{,}46 \cdot 10^{-2}$, cao hơn cả bước hằng; với $\gamma = \mu \eta_{0}$
theo lý thuyết cho hàm lồi mạnh, tức nhỏ hơn 30 lần, sai số xuống
$6{,}94 \cdot 10^{-6}$. Ngân sách cũng quyết định kết luận: ở 40 epoch bước hằng
vẫn đang dẫn, và chỉ từ khoảng 100 epoch trở đi thứ hạng mới đảo.

Bước hằng còn khác ba quy tắc kia ở mức phân tán giữa các lần chạy. Sai số cuối
của nó trải từ $2{,}59 \cdot 10^{-4}$ tới $9{,}64 \cdot 10^{-4}$ qua 5 seed, tức
chênh $3{,}7$ lần, trong khi ba quy tắc giảm dần cho kết quả gần trùng nhau giữa
các seed. Điểm dừng của bước hằng nằm trong một lân cận ngẫu nhiên nên phụ thuộc
vào lô cuối cùng, còn quy tắc giảm dần đã co dãy lặp về gần một điểm tất định.
Hình~\ref{fig:band} vẽ trung vị kèm dải min-max qua 5 seed, còn
hình~\ref{fig:schedule} vẽ riêng seed 0 theo hai trục đo.

\resultfig{batch-schedule_band}{Bốn quy tắc chọn bước, trung vị và dải min-max
  qua 5 seed.}{fig:band}

\resultpair{batch-schedule}{Bốn quy tắc chọn bước cho SGD, seed
  0.}{fig:schedule}

% ===========================================================================
\chapter{Thông tin bậc hai: đổi giá mỗi vòng lấy số vòng lặp}\label{ch:newton}
% ===========================================================================

\section{Một vòng lặp, sai số bằng không}

Phương pháp Newton giải hệ $\nabla^{2} f(w_k)\, p_k = -\nabla f(w_k)$ rồi cập
nhật $w_{k+1} = w_k + t\, p_k$. Hệ được giải bằng phân rã Cholesky
\cite{nocedal2006numerical} chứ không bằng cách tính nghịch đảo tường minh, vì
nghịch đảo vừa đắt hơn vừa kém ổn định số học.

Trên hàm mục tiêu \eqref{eq:objective}, Newton với $t = 1$ hội tụ sau đúng một
vòng lặp và cho $f - \fstar$ bằng $0$ đúng nghĩa trong \texttt{float64}. Kết quả
này suy ra thẳng từ \eqref{eq:grad}, vì Hessian là hằng số nên xấp xỉ bậc hai
trùng khít hàm thật, và với $w_{0} = 0$ thì
\begin{equation}\label{eq:newtonstep}
  w_{1} = w_{0} - \left( \nabla^{2} f \right)^{-1} \nabla f(w_{0})
        = \left( \tfrac{1}{n} X\trans X + \lambda I \right)^{-1} \tfrac{1}{n} X\trans y = \wstar,
\end{equation}
tức bước Newton \eqref{eq:newtonstep} chính là nghiệm đóng \eqref{eq:closed}.
Hình~\ref{fig:newton} vẽ cả bốn biến thể. Backtracking cũng vì thế
không có việc gì để làm: điều kiện giảm đủ thỏa ngay ở lần thử thứ nhất, đo được
đúng một lần đánh giá hàm mục tiêu cho toàn bộ lần chạy. Phần backtracking cho
Newton vì thế chỉ có nội dung khi hàm mục tiêu không còn bậc hai, chẳng
hạn khi thay mất mát bình phương bằng mất mát Huber.

\resultpair{newton-damping}{Newton với bước đầy đủ và bước damped, hai cách xử lý
  Hessian.}{fig:newton}

\section{Giá của một vòng lặp}

Điều đáng so sánh với Newton không phải số vòng lặp mà là chi phí của vòng lặp
duy nhất ấy. Bảng~\ref{tab:newtoncost} tách chi phí đó thành ba phần, đo ở quy mô
$n = \num{1200000}$.

\begin{table}[tbp]
  \centering
  \caption{Chi phí từng phần của một bước Newton, $n = \num{1200000}$,
    $d = 116$.}
  \label{tab:newtoncost}
  \begin{tabular}{lrl}
    \toprule
    Thành phần & Thời gian (ms) & Bậc chi phí \\
    \midrule
    Dựng $\frac{1}{n} X\trans X$      & 163{,}2 & $\bigO(nd^{2})$ \\
    Phân rã Cholesky $116 \times 116$ & 0{,}2   & $\bigO(d^{3})$ \\
    Một lần tính gradient             & 87{,}3  & $\bigO(nd)$ \\
    \bottomrule
  \end{tabular}
\end{table}

Một bước Newton đầy đủ tốn khoảng $250$ ms, trong đó phần $\bigO(d^{3})$ chỉ
chiếm $0{,}2$ ms, tức hoàn toàn không đáng kể. Với $d = 116$ và $n$ cỡ triệu, chi
phí dựng Hessian chỉ đắt hơn một lần tính gradient khoảng hai lần chứ không phải
$d$ lần như dự đoán thô ban đầu, vì phép nhân ma trận với ma trận thuộc nhóm BLAS
bậc ba, đọc mỗi phần tử một lần rồi dùng lại trong bộ nhớ đệm, trong khi tính
gradient phải quét toàn bộ ma trận hai lượt nên băng thông bộ nhớ chặn nó lại.
Kết luận đảo chiều khi $d$ lên tới vài nghìn, vì chi phí Hessian tăng theo $d^{2}$
còn chi phí gradient chỉ tăng theo $d$.

Một đính chính về cách đo cần nêu rõ. Biến thể dùng lại phân rã Cholesky đo được
$0{,}16$ giây, nhưng con số đó bỏ sót phần dựng Hessian, vì thủ tục làm nóng chạy
trước khi bấm đồng hồ đã tính sẵn ma trận đó. Biến thể dựng lại Hessian mỗi vòng
đo được $0{,}40$ giây, cao hơn $250$ ms vì vòng lặp gọi thêm một lần đề xuất bước
ở vòng thứ hai chỉ để kiểm tra điều kiện dừng, và với Newton thì lần gọi đó kéo
theo cả việc dựng lại Hessian. Con số dùng để so sánh trong báo cáo là $250$ ms.

% ===========================================================================
\chapter{So sánh tổng hợp trên hai trục}\label{ch:headline}
% ===========================================================================

\section{Cấu hình tốt nhất của từng phương pháp}

Bảng~\ref{tab:headline} tập hợp cấu hình tốt nhất tìm được ở các chương trước, đo
lại ở cả hai quy mô mẫu với ba lần chạy độc lập mỗi cấu hình và lấy trung vị.

\begin{table}[tbp]
  \centering
  \caption{Cấu hình tốt nhất của từng phương pháp ở hai quy mô. Cột số vòng lặp
    là số vòng chạy hết, cột thời gian là thời gian tương ứng.}
  \label{tab:headline}
  \begin{tabular}{lrrrr}
    \toprule
    & \multicolumn{2}{c}{$n = \num{200000}$} & \multicolumn{2}{c}{$n = \num{1200000}$} \\
    \cmidrule(lr){2-3}\cmidrule(lr){4-5}
    Cấu hình & Vòng lặp & Giây & Vòng lặp & Giây \\
    \midrule
    Newton, Hessian dùng lại        & 1     & 0{,}03  & 1     & 0{,}16  \\
    AGD, $\beta$ theo $t$ và $\mu$  & 286   & 3{,}92  & 283   & 23{,}49 \\
    GD, backtracking                & 561   & 13{,}86 & 514   & 71{,}99 \\
    GD, $t = 1{,}9/L$               & 2022  & 27{,}52 & 1959  & 158{,}67 \\
    SGD, $B = 2048$                 & 3880  & 2{,}02  & 23400 & 13{,}17 \\
    \bottomrule
  \end{tabular}
\end{table}

Thứ hạng theo hai trục đo giống nhau ở bài toán này, khác với tình huống mà đề
bài dự đoán. Newton đứng đầu cả về số vòng lặp lẫn thời gian, và khoảng cách
theo thời gian với gradient descent là gần một nghìn lần ở quy mô toàn phần. Lý
do là giá mỗi vòng của bốn phương pháp chênh nhau ít hơn số vòng lặp chênh nhau,
vì một vòng Newton đắt hơn một vòng gradient descent khoảng ba lần, trong khi
gradient descent cần nhiều hơn hai nghìn lần số vòng. Thứ hạng sẽ đảo khi $d$
lớn, vì lúc đó chi phí Hessian tăng theo $d^{2}$.

SGD là ngoại lệ và phải đọc riêng. Nó nhanh nhất trên mỗi đơn vị thời gian, đạt
sai số $6{,}0 \cdot 10^{-4}$ chỉ sau $13{,}17$ giây ở quy mô toàn phần, nhưng
không bao giờ tới ngưỡng $10^{-6}$ vì sàn nhiễu mô tả ở chương~\ref{ch:sgd}. Nếu
bài toán chỉ đòi độ chính xác cỡ $10^{-3}$ thì SGD thắng tuyệt đối; ở ngưỡng chặt
hơn thì nó không phải một lựa chọn. Hình~\ref{fig:headline} và
hình~\ref{fig:headlinefull} vẽ hai quy mô.

\resultpair{headline}{Cấu hình tốt nhất của từng phương pháp, $n =
  \num{200000}$.}{fig:headline}

\resultpair{headline-full}{Cấu hình tốt nhất của từng phương pháp, $n =
  \num{1200000}$.}{fig:headlinefull}

\section{Ảnh hưởng của kích thước mẫu}

So sánh hai nửa của bảng~\ref{tab:headline} cho thấy số vòng lặp gần như không
đổi khi $n$ tăng sáu lần: gradient descent đi từ 2022 xuống 1959 vòng, gradient
tăng tốc từ 286 xuống 283. Thời gian thì tăng từ $5{,}2$ tới $6{,}5$ lần, tức
tuyến tính theo $n$ đúng như mô hình chi phí $\bigO(nd)$ mỗi vòng.

Hai quan sát này là hệ quả trực tiếp của hệ số $\frac{1}{n}$ trong
\eqref{eq:objective}. Nó giữ $\kappa$ gần như cố định giữa hai quy mô, $267{,}5$
so với $266{,}2$, mà số vòng lặp của mọi phương pháp bậc một lại chỉ phụ thuộc
$\kappa$ chứ không phụ thuộc $n$. Bỏ hệ số đó đi thì $L$ tăng tuyến tính theo $n$
và toàn bộ lưới tham số phải quét lại ở mỗi quy mô. SGD là ngoại lệ vì ngân sách
của nó đặt theo số epoch, mà một epoch ở quy mô lớn chứa nhiều lô hơn.

\section{Đối chiếu với tốc độ lý thuyết}

Lý thuyết cho hàm bậc hai lồi mạnh nói số vòng lặp của gradient descent tỷ lệ với
$\kappa$, còn của gradient tăng tốc tỷ lệ với $\sqrt{\kappa}$
\cite{nesterov2018lectures}. Bảng~\ref{tab:theory} đối chiếu bằng cách đo trên
cùng bộ dữ liệu với các giá trị $\lambda$ khác nhau, mỗi lần dùng bước tính theo
$L$ và $\mu$ của chính bài toán đó.

\begin{table}[tbp]
  \centering
  \caption{Tỷ lệ số vòng lặp giữa các cặp $\kappa$ liền kề, so với tỷ lệ mà lý
    thuyết dự đoán.}
  \label{tab:theory}
  \begin{tabular}{lrrrr}
    \toprule
    $\kappa$ giảm & Tỷ lệ $\kappa$ & GD quan sát & AGD quan sát & $\sqrt{\kappa}$ dự đoán \\
    \midrule
    730 xuống 268 & 2{,}73 & 2{,}73 & 1{,}57 & 1{,}65 \\
    268 xuống 90  & 2{,}98 & 2{,}98 & 1{,}56 & 1{,}73 \\
    90 xuống 10   & 8{,}91 & 8{,}86 & 2{,}82 & 2{,}99 \\
    \bottomrule
  \end{tabular}
\end{table}

Cột gradient descent khớp cận lý thuyết tới hai chữ số ở cả ba cặp. Cột gradient
tăng tốc bám sát $\sqrt{\kappa}$ nhưng lệch chừng năm tới mười phần trăm, và lệch
theo hướng tốt hơn dự đoán ở cặp cuối. Hằng số nhân trước lũy thừa trong cận tăng
tốc phụ thuộc $\kappa$, nên tỷ lệ giữa hai lần chạy không rút gọn sạch như ở cận
của gradient descent.

% ===========================================================================
\chapter{So sánh với scikit-learn}\label{ch:library}
% ===========================================================================

\section{Quy đổi hàm mục tiêu}

Điểm dễ sai nhất khi so sánh với thư viện là mỗi ước lượng chuẩn hóa hàm mục tiêu
theo một cách khác nhau. Lớp \texttt{Ridge} của scikit-learn
\cite{pedregosa2011scikit} cực tiểu hóa $\norm{Xw-y}^{2} + \alpha \norm{w}^{2}$,
tức bỏ cả hệ số $\frac{1}{n}$ lẫn hệ số $\frac{1}{2}$ của \eqref{eq:objective};
nhân \eqref{eq:objective} với $2n$ cho quan hệ
\begin{equation}\label{eq:alpha}
  \alpha_{\text{Ridge}} = \lambda n .
\end{equation}
Lớp \texttt{SGDRegressor} lại cực tiểu hóa
$\frac{1}{n}\sum_i \frac{1}{2}(x_i\trans w - y_i)^{2} + \frac{\alpha}{2}\norm{w}^{2}$,
tức trùng đúng với \eqref{eq:objective}, nên với nó $\alpha = \lambda$. Hai ước
lượng trong cùng một thư viện cần hai hằng số khác nhau, và đặt sai không sinh ra
lỗi nào: cả hai bên vẫn giải sạch, chỉ là giải hai bài toán khác nhau.

Quy đổi \eqref{eq:alpha} được kiểm chứng bằng số chứ không bằng suy luận trên
giấy. Lấy nghiệm $\hat{w}$ do \texttt{Ridge} trả về với
$\alpha = \lambda n = 3{,}795 \cdot 10^{4}$ rồi so với nghiệm đóng cho
$\norm{\hat{w} - \wstar} / \norm{\wstar} = 2{,}35 \cdot 10^{-15}$, tức ở mức sai
số máy. Phép kiểm này chạy tự động trước mọi so sánh và làm dừng chương trình nếu
không đạt.

\section{Kết quả}

Bảng~\ref{tab:library} ghi kết quả ở quy mô toàn phần.

\begin{table}[tbp]
  \centering
  \caption{So sánh với scikit-learn ở quy mô $n = \num{1200000}$, trung vị của
    ba lần chạy. Dòng đầu là nghiệm đóng của nhóm.}
  \label{tab:library}
  \begin{tabular}{lrrr}
    \toprule
    Phương pháp & Thời gian (s) & $f - \fstar$ & RMSE test \\
    \midrule
    Newton của nhóm            & 0{,}25  & $0$                    & 3{,}4646 \\
    \texttt{Ridge(cholesky)}   & 0{,}273 & $4{,}15 \cdot 10^{-30}$ & 3{,}4646 \\
    \texttt{Ridge(lsqr)}       & 2{,}03  & $1{,}13 \cdot 10^{-5}$  & 3{,}4647 \\
    \texttt{Ridge(sparse\_cg)} & 2{,}20  & $1{,}99 \cdot 10^{-7}$  & 3{,}4646 \\
    \texttt{Ridge(sag)}        & 45{,}44 & $3{,}96 \cdot 10^{-7}$  & 3{,}4647 \\
    \texttt{SGDRegressor}      & 3{,}18  & $6{,}21 \cdot 10^{20}$  & $3{,}3 \cdot 10^{10}$ \\
    \texttt{LinearRegression}  & 3{,}08  & $1{,}75 \cdot 10^{-2}$  & 3{,}4614 \\
    \bottomrule
  \end{tabular}
\end{table}

Mã tự viết và thư viện làm cùng một việc và mất cùng một thời gian: $0{,}25$ giây
so với $0{,}273$ giây, với sai số của cả hai ở mức sai số máy. Hai con số sát nhau
không nói bên nào nhanh hơn, mà là bằng chứng cả hai cài đúng, vì
\texttt{Ridge} với solver mặc định cũng chính là nghiệm đóng giải qua Cholesky.
Mục đích của việc tự cài đặt không phải chạy nhanh hơn thư viện, mà là hiểu cơ
chế hội tụ và biết cách chọn tham số.

Ba solver lặp của \texttt{Ridge} đổi độ chính xác lấy thời gian theo những tỷ lệ
rất khác nhau. Solver \texttt{lsqr} dừng ở sai số $1{,}13 \cdot 10^{-5}$ sau $2{,}03$
giây, còn \texttt{sag} tốn $45{,}44$ giây để đạt $3{,}96 \cdot 10^{-7}$, tức chậm
hơn nghiệm trực tiếp 166 lần mà vẫn kém chính xác hơn 23 bậc. Chúng chỉ có lợi
khi $d$ lớn tới mức chi phí $\bigO(d^{3})$ hoặc bộ nhớ cho $X\trans X$ không chấp
nhận được, điều không xảy ra với $d = 116$.

\texttt{SGDRegressor} với tham số mặc định cho sai số $6{,}21 \cdot 10^{20}$ và
RMSE $3{,}3 \cdot 10^{10}$, tức hoàn toàn không dùng được, sau khi dừng ở vòng
lặp thứ sáu theo tiêu chí dừng mặc định. Kết quả này là phát biểu về giá trị mặc
định chứ không phải về thuật toán: SGD tự cài với bước đặt theo $1/L_{B}$ đạt sai
số $6{,}0 \cdot 10^{-4}$ trên cùng bài toán. Chỉnh tham số cho
\texttt{SGDRegressor} sẽ cải thiện, nhưng mặc định của thư viện không đặt bước
theo $L$ của bài toán cụ thể.

Dòng \texttt{LinearRegression} ứng với $\lambda = 0$ nên giải một bài toán khác,
và sai số $1{,}75 \cdot 10^{-2}$ so với $\fstar$ phản ánh đúng điều đó. Đáng chú
ý là RMSE trên tập kiểm tra của nó, $3{,}4614$, lại nhỉnh hơn nghiệm Ridge một
chút, $3{,}4646$. Chương~\ref{ch:lambda} bàn về khoản chênh này. Hình~\ref{fig:library} đặt các kết
quả thư viện cạnh đường hội tụ của mã tự viết.

\resultpair{library-full}{So sánh mã tự viết với scikit-learn ở quy mô
  $n = \num{1200000}$. Các kết quả thư viện là điểm đơn vì solver không cho lịch
  sử theo vòng lặp.}{fig:library}

% ===========================================================================
\chapter{Ảnh hưởng của hệ số hiệu chỉnh}\label{ch:lambda}
% ===========================================================================

Hệ số $\lambda$ có hai vai trò tách biệt. Vai trò thống kê là chống quá khớp, và
vai trò tối ưu hóa là đặt sàn cho phổ của Hessian: theo \eqref{eq:constants} thì
$\mu = \lambda_{\min} + \lambda$, nên tăng $\lambda$ làm tăng $\mu$, giảm $\kappa$
và do đó tăng tốc độ hội tụ. Bảng~\ref{tab:lambda} đo cả hai vai trò cùng lúc.

\begin{table}[tbp]
  \centering
  \caption{Ảnh hưởng của $\lambda$ lên số điều kiện, tốc độ hội tụ và chất lượng
    dự báo, $n = \num{200000}$. Cột vòng lặp là số vòng cần để đạt
    $f - \fstar < 10^{-6}$.}
  \label{tab:lambda}
  \begin{tabular}{rrrrrr}
    \toprule
    $\lambda$ & $\mu$ & $\kappa$ & GD & AGD & RMSE test \\
    \midrule
    0{,}001    & 0{,}00345  & 2633{,}0 & không đạt & 209 & 3{,}4537 \\
    0{,}01     & 0{,}01245  & 730{,}4  & 2091      & 118 & 3{,}4541 \\
    0{,}031623 & 0{,}03407  & 267{,}5  & 766       & 75  & 3{,}4563 \\
    0{,}1      & 0{,}10245  & 89{,}6   & 257       & 48  & 3{,}4690 \\
    1          & 1{,}00245  & 10{,}1   & 29        & 17  & 3{,}7270 \\
    10         & 10{,}00245 & 1{,}9    & 5         & 6   & 4{,}4578 \\
    \bottomrule
  \end{tabular}
\end{table}

Tồn tại một khoảng $\lambda$ mua được điều kiện gần như miễn phí. Đi từ
$\lambda = 0{,}001$ lên $\lambda = 0{,}1$ làm RMSE trên tập kiểm tra xấu đi từ
$3{,}4537$ lên $3{,}4690$, tức $0{,}44$ phần trăm, trong khi $\kappa$ giảm 29 lần
và gradient descent đi từ chỗ không hội tụ nổi sau \num{4000} vòng xuống còn 257
vòng. Đường cong cross-validation phẳng trên khoảng đó, vì $\lambda$ trong vùng này chưa
đủ lớn để làm lệch nghiệm một cách đáng kể nhưng đã đủ lớn để nâng hẳn sàn của
phổ. Ra ngoài khoảng đó thì cái giá hiện ngay, với $\lambda = 1$ cho
RMSE $3{,}7270$ và $\lambda = 10$ cho $4{,}4578$, tức xấu đi 29 phần trăm.

Giá trị $\lambda = 0{,}031623$ mà quy tắc một sai số chuẩn chọn ở
chương~\ref{ch:step} nằm đúng trong khoảng đó, với $\kappa = 267{,}5$ và RMSE
$3{,}4563$. Cũng cần ghi nhận rằng nghiệm không hiệu chỉnh cho RMSE $3{,}4614$,
tức nhỉnh hơn nghiệm Ridge được chọn; trên bộ dữ liệu này với $n$ lớn hơn $d$ hơn
một nghìn lần, hiệu chỉnh gần như không mua thêm được gì về mặt thống kê, và toàn
bộ lợi ích của nó nằm ở phía tối ưu hóa. Hình~\ref{fig:lambda} vẽ các đường hội tụ
tương ứng, còn hình~\ref{fig:tradeoff} đặt $\kappa$ và RMSE lên cùng một khung.

\resultpair{lambda-kappa}{Gradient descent với các giá trị $\lambda$ khác nhau,
  mỗi lần dùng bước $2/(L+\mu)$ của chính bài toán đó.}{fig:lambda}

\resultfig{lambda-kappa_tradeoff}{Số điều kiện và sai số dự báo vẽ cùng theo
  $\lambda$ trên hai trục tung.}{fig:tradeoff}

% ===========================================================================
\chapter{Kết luận}
% ===========================================================================

Báo cáo so sánh bốn thuật toán tối ưu hóa trên cùng một hàm mất mát Ridge dựng từ
dữ liệu Lending Club, với $\kappa = 267{,}5$ và $d = 116$. Ở cấu hình tốt nhất
của từng phương pháp, tổng hợp ở chương~\ref{ch:headline}, thứ hạng theo cả hai
trục đo là Newton, gradient tăng tốc,
gradient descent với backtracking, rồi gradient descent với bước cố định; SGD
nhanh nhất trên mỗi đơn vị thời gian nhưng dừng ở sai số $6{,}0 \cdot 10^{-4}$.

Ba điều rút ra về việc chọn tham số. Với backtracking, $\rho$ quan trọng hơn $c$
rất nhiều, vì nó quyết định số lần thử mỗi vòng. Với SGD, bước phải đặt theo
$1/L_{B}$ của kích thước lô chứ không theo $1/L$, và tốc độ giảm phải neo vào
$\mu$ chứ không vào số vòng lặp mỗi epoch. Với gradient tăng tốc, $\beta$ phải
tính lại theo bước mà line search thực sự chọn, như chương~\ref{ch:momentum} chỉ
ra.

Ba chỗ thực nghiệm lệch khỏi phát biểu quen thuộc của lý thuyết, và cả ba đều
lệch vì cùng một lý do: cận lý thuyết là cận cho trường hợp xấu nhất trên một
lớp bài toán, còn thực nghiệm chạy trên một bài toán cụ thể, một điểm khởi tạo
cụ thể và một ngân sách hữu hạn. Bước $1{,}9/L$ nhanh gấp $2{,}3$ lần bước tối ưu
$2/(L+\mu)$, vì sai số ban đầu tại $w_{0} = 0$ dồn $80{,}5$ phần trăm vào các
hướng trị riêng lớn. Backtracking thắng bước cố định trên cả hai trục, vì nó so
với thương Rayleigh tại chỗ chứ không với $L$ toàn cục. Và quy tắc giảm bước của
SGD chỉ thắng bước hằng khi tốc độ giảm đặt đúng, với hai giá trị hợp lý cho hai
kết luận ngược nhau.

Kết luận cuối cùng về phạm vi áp dụng. Với $d = 116$ và $n$ cỡ triệu, phương pháp
bậc hai thắng áp đảo, vì chi phí dựng Hessian là $\bigO(nd^{2})$ với $d^{2}$ nhỏ,
còn phần $\bigO(d^{3})$ chỉ chiếm $0{,}2$ ms trong tổng $250$ ms. Các phương pháp
lặp bậc một chỉ thể hiện lợi thế khi $d$ lớn tới mức không dựng nổi $X\trans X$,
hoặc khi dữ liệu không nạp hết vào bộ nhớ. Đối chiếu với scikit-learn ở chương~\ref{ch:library} cho thấy nghiệm đóng tự cài
và solver mặc định của thư viện mất cùng một thời gian, vì cả hai làm cùng một
phép tính. Tính $L$ và $\mu$ trước khi chọn bước tốn dưới một giây trên bài toán
này, trong khi quét năm bước theo lối dò tay
tốn 195 giây và vẫn không tìm ra bước gần tối ưu.

\printbibliography[heading=bibintoc, title={Tài liệu tham khảo}]

\end{document}
